% Generated draft: requires mathematical and source review.
\documentclass[11pt]{article}
\usepackage[margin=25mm]{geometry}
\usepackage{fontspec}
\setmainfont{Latin Modern Roman}
\usepackage{amsmath,amssymb,mathtools,unicode-math}
\setmathfont{Latin Modern Math}
\usepackage{enumitem,needspace,xurl,hyperref}
\setlength{\parindent}{0pt}
\setlength{\parskip}{5pt}
\setlength{\emergencystretch}{3em}
\newcommand{\N}{\mathbb{N}}
\newcommand{\Z}{\mathbb{Z}}
\newcommand{\Q}{\mathbb{Q}}
\newcommand{\R}{\mathbb{R}}
\newcommand{\C}{\mathbb{C}}
\newcommand{\F}{\mathbb{F}}
\newcommand{\A}{\mathbb{A}}
\newcommand{\PP}{\mathbb{P}}
\newcommand{\E}{\mathbb{E}}
\newcommand{\eps}{\varepsilon}
\newcommand{\dd}{\,\mathrm{d}}
\begin{document}
\section*{UnsolvedMath ID: 2400005}
\textbf{Source problem code:} catalogue entry
\subsection{Definitions and mathematical statement}
The intended root-finding setting is a local iteration near a zero $\alpha$ of a scalar function $f$. Write $x_{k+1}=\Phi_f(x_k)$ and $e_k=x_k-\alpha$. A standard exact local order $p$ is expressed by
\[
\lim_{k\to\infty}\frac{|e_{k+1}|}{|e_k|^p}=C,
\qquad 0<C<\infty,
\]
for sufficiently close initial points and the intended generic class of functions. ``Without memory'' means that each step uses only the current iterate and information acquired during that step, without retaining earlier-step evaluation data; intermediate multipoint evaluations within the step are allowed.

The Kung--Traub conjecture asks whether every admissible method using $n\geq1$ evaluations of $f$ or its derivatives per step has order
\[
p\leq2^{n-1}.
\]
The evaluation model is part of the assertion: which derivative orders can be queried, how each query is charged, which operations may combine the data, and what counts as an evaluation must be specified. The usual conjecture concerns a method's order for a general sufficiently regular function near a simple zero, not exceptional exact termination for a special input. The source does not state a function class, root-multiplicity assumption, generic-order convention, or a full computational model, so these standard contextual choices are not silently promoted to verified source hypotheses.

The supplied research summary attributes the conjecture to Kung--Traub (1974), reports it established for small $n$ and particular families, and leaves the general statement unresolved. It mentions Newton and optimal multipoint methods as low-evaluation examples, but supplies no exact low-$n$ theorem or cost convention. Under the usual separate counting of $f$ and $f'$ queries, Newton uses two evaluations per step. The source asks for the upper bound itself and does not assert that every method attains it. The summary reports no known general proof or counterexample; that is a supplied literature assessment, not a new verification.
\subsection{Short English statement}
Kung--Traub predicts that a memoryless scalar root-finding step with $n$ function or derivative evaluations cannot have generic order above $2^{n-1}$. The source reports the general bound open, but its evaluation model, function class, and local-order convention need to be supplied for a precise theorem.
\subsection{Sources}
\begin{itemize}
\item Local UnsolvedMath registry, permanent ID 2400005, exact source code catalogue entry, ``Kung--Traub conjecture''; status open, classification OPEN-TRIAGE. The source provenance assigned default L3; the later research recommendation and current metadata give L4.
\item Wikipedia, \emph{List of unsolved problems in mathematics}, pinned revision 1366636767, analysis item 5, \url{https://en.wikipedia.org/w/index.php?title=List_of_unsolved_problems_in_mathematics&oldid=1366636767}, accessed July 29, 2026. The record identifies adaptation from this pinned revision under CC BY-SA 4.0, with attribution, and says primary-source/current-status release review is still required.
\item Kung and Traub (1974), cited by the supplied research summary without a full title, publication details, or precise model statement. The summary reports low-evaluation and special-family progress but provides no exact theorem references, and describes the general conjecture as unresolved.
\end{itemize}
\end{document}
